\unchapter{主要符号与缩略语}

\noindent 本表列出全文共用概念。各研究章中的时隙、任务和决策步骤按其实际过程定义，具体下标及局部符号在对应模型中说明。

\begin{longtable}{@{}p{31mm}p{111mm}@{}}
\toprule
符号或缩略语&含义\\
\midrule
\endfirsthead
\toprule 符号或缩略语&含义\\\midrule
\endhead
\bottomrule\endfoot
$\mathcal G(t)$&时变通信图\\
$\mathcal V,\mathcal E(t)$&节点集合与可用链路集合\\
$\mathcal N_i(t)$&节点 $i$ 的当前邻居集合\\
$Q_i(t)$&节点数据队列\\
$E_i(t)$&节点剩余能量\\
$F_i(t),Z_i(k)$&路由时隙或语义任务尺度上的虚拟能量队列\\
$r_{ij}(t),P_i(t)$&链路速率与节点发送功率\\
$\bm b^{\rm sem}$&任务类别上的语义信念分布\\
$b^{\rm rec}(S)$&实际接收集合上的状态概率\\
$\mathcal M_k$&任务 $k$ 的候选语义证据集合\\
$S_{k,t}$&任务内实际已经接收的唯一证据集合\\
$D_k$&任务相对截止时间\\
$p_m^{\rm del}$&候选消息的预测送达概率\\
$q_k(S)$&经一致性修正的集合任务效用\\
$\pi_\theta$&参数为 $\theta$ 的通信选择策略\\
UAV&Unmanned Aerial Vehicle，无人机\\
FANET&Flying Ad Hoc Network，飞行自组网\\
VLM&Vision-Language Model，视觉语言模型\\
VQA&Visual Question Answering，视觉问答\\
PPO&Proximal Policy Optimization，近端策略优化\\
GAE&Generalized Advantage Estimation，广义优势估计\\
ACK&Acknowledgment，确认反馈\\
ARQ&Automatic Repeat Request，自动重传请求\\
HARQ&Hybrid Automatic Repeat Request，混合自动重传请求\\
POMDP&Partially Observable Markov Decision Process，部分可观测马尔可夫决策过程\\
LEADR&Lyapunov-based Energy-Aware Decentralized Routing，基于李雅普诺夫优化的能量感知分布式路由\\
EviComRL&基于强化学习的链路感知选择性语义通信方法\\
C-BFC&Consistency-aware Belief Feedback Control，一致性感知的信念反馈控制\\
SD-SAC&Stable Discrete Soft Actor-Critic，稳定离散软演员--评论家方法\\
TP-DQN&第~\ref{rq3:chapter} 章适配的两阶段性能预测与深度 Q 网络控制器\\
\end{longtable}
